\documentclass[10pt,a4paper]{amsart}
\usepackage[margin=19mm]{geometry}
\usepackage{amsmath,amssymb,mathtools}
\usepackage[scr=rsfs,cal=euler]{mathalfa}
\usepackage{xfrac}
\usepackage[unicode,hidelinks,bookmarks=false]{hyperref}
\newcommand{\ie}{i.e.\ }
\newcommand{\cf}{cf.\ }
\newcommand{\resp}{respectively\ }
\newcommand{\Subsection}{\subsection}
\providecommand{\nobreakdash}{\nobreak\hbox{-}}
\setlength{\emergencystretch}{2em}
\multlinegap0pt
\usepackage{amssymb}
\usepackage{mathtools}
\usepackage{xcolor}
\usepackage[normalem]{ulem}
\usepackage[shortlabels]{enumitem}
\newtheorem{theorem}{Theorem}
\title{Asymptotic behavior of Eckhoff's method for convergence acceleration of Dirac eigenfunction expansions}
\author{R.H. Barkhudaryan, G.G. Gevorkyan, L.D. Poghosyan}
\date{Source: arXiv:2609.00895v1 (CC BY 4.0)}
\begin{document}
\maketitle

\noindent\emph{Source and licence.} Asymptotic behavior of Eckhoff's method for convergence acceleration of Dirac eigenfunction expansions. Version arXiv:2609.00895v1 (CC BY 4.0), distributed by arXiv under the Creative Commons Attribution 4.0 International licence, http://creativecommons.org/licenses/by/4.0/, as recorded in the arXiv licence metadata for this exact version and shown on the abstract page https://arxiv.org/abs/2609.00895v1. Full licence notice: \texttt{licenses/arxiv-2609.00895-CC-BY-4.0.txt}.\par\smallskip

\noindent\textit{Editorial scope.} Counted unit: the article's theorem thm:L2 (source0.tex lines 2274--2287). The article's complete original proof of it is delivered with it (source0.tex lines 2289--2574, one \texttt{proof} environment of 286 lines). No mathematics is written, repaired or re-derived: the statement and the proof are the article's own lines, in the article's own notation, and no retained line is substituted or re-flowed.  Retained uncounted as the source-local context the proof needs: lines 208--218; lines 490--493; lines 522--527; lines 1162--1166; lines 1285--1287; lines 1307--1315; lines 1441--1443; lines 1503--1507; lines 1687--1705; lines 2080--2084; lines 2098--2135.  Nothing was omitted from any retained span, so the package carries no omission record.  No retained line contains a citation, so the wrapper carries no bibliography and no bibliography entry.  No retained line contains a figure environment, an image-inclusion command or drawing code, so no image is delivered and the package declares no asset dependency.  Editorial formatting only, in addition: an AMS \texttt{amsart} preamble that restates the article's own declarations and macros used by the retained lines, namely the environments \texttt{theorem} on separate counters, together with the standard packages those lines need.  The retained environments print in this document as Theorem 1 in order of appearance, whereas the article prints the same items with the numbering of its own full document.  The complete native source of the article as distributed in the arXiv e-print for this exact version is bundled unchanged as \texttt{sources/209068/source0.tex}.
\begin{equation}\label{eq:lam-asymp}
\lambda_n
=
\frac{\pi n}{2}
-
\frac{\widetilde\theta}{2}
+
O\!\left(\frac{1}{n}\right),
\qquad
n\to\pm\infty,
\end{equation}

\begin{equation}\label{eq:Fn-asymp}
F_n=x_n^{q+1}\bigl[\kappa_n^{(+)}g_q^+-\kappa_n^{(-)}g_q^-\bigr]+o(|n|^{-q-1}),
\qquad |n|\to\infty.
\end{equation}

\begin{equation}\label{eq:AB}
\begin{aligned}
A&:=g_q^+ = f_{q,1}(1)\cos\beta - f_{q,2}(1)\sin\beta,\\
B&:=-g_q^- = f_{q,2}(-1)\sin\alpha - f_{q,1}(-1)\cos\alpha.
\end{aligned}
\end{equation}

\begin{equation}\label{eq:ns-range}
\alpha_0 N\leq |n_s|\leq N,
\qquad
s=1,\dots,2q,
\end{equation}

\begin{equation}\label{eq:Phi-def}
\Phi^\pm(x):=\Delta^\pm(x)-g_q^\pm x^q\qquad(\deg\le q),
\end{equation}

\begin{equation}\label{eq:kappa-leading}
\kappa_n^{(+)}
=
(-1)^nK_0+O\!\left(\frac{1}{n}\right),
\qquad
-\kappa_n^{(-)}
=
-K_0+O\!\left(\frac{1}{n}\right),
\end{equation}

\begin{equation}\label{eq:parity-balance}
q_e=q_o=q.
\end{equation}

\begin{equation}\label{eq:Pi-EO}
\Pi_E(x):=\prod_{s\in E}(x-x_{n_s}),
\qquad
\Pi_O(x):=\prod_{s\in O}(x-x_{n_s}).
\end{equation}

\begin{equation}\label{eq:UV-explicit}
\begin{aligned}
U^+(x)
&=
\frac{1}{2}\bigl[\Pi_E(x)+\Pi_O(x)\bigr],
&
U^-(x)
&=
-\frac{1}{2}\bigl[\Pi_E(x)-\Pi_O(x)\bigr],
\\[2pt]
V^+(x)
&=
-\frac{1}{2}\bigl[\Pi_E(x)-\Pi_O(x)\bigr],
&
V^-(x)
&=
\frac{1}{2}\bigl[\Pi_E(x)+\Pi_O(x)\bigr].
\end{aligned}
\end{equation}

\begin{equation}\label{eq:Phi-decomp-2}
\Phi^\pm(x)
=
-g_q^+U^\pm(x)-g_q^-V^\pm(x)+\rho^\pm(x),
\end{equation}

\begin{theorem}\label{thm:gjumps}
Suppose
\[
f\in C^q([-1,1];\mathbb R^2),
\qquad
p,r\in C^q[-1,1],
\]
with $f^{(q)}$ absolutely continuous, where $q\geq1$. Assume that the
indices $n_s=n_s(N)$ satisfy \eqref{eq:ns-range} with
\[
\lim_{N\to\infty}\frac{n_s}{N}=c_s\neq0,
\qquad
s=1,\dots,2q,
\]
that the strict-balance condition \eqref{eq:parity-balance} holds, and
that the limits $\{c_s:s\in E\}$ are pairwise distinct, as are
$\{c_s:s\in O\}$. Then
\begin{equation}\label{eq:tilde-g-asymp}
\widetilde g_k^\pm
=
g_k^\pm
-g_q^+U_k^\pm
-g_q^-V_k^\pm
+o(N^{k-q}),
\qquad
k=0,\dots,q-1,
\end{equation}
where $U_k^\pm$ and $V_k^\pm$ denote the coefficients of $x^k$ in
the polynomials $U^\pm$ and $V^\pm$ given by
\eqref{eq:UV-explicit}. In particular,
\[
\widetilde g_k^\pm-g_k^\pm
=
O(N^{k-q}),
\qquad
k=0,\dots,q-1.
\]
\end{theorem}

\begin{theorem}\label{thm:L2}
Under the hypotheses of Theorem~\ref{thm:gjumps} (in particular, the
strict parity balance $q_e=q_o=q$),
\begin{equation}\label{eq:L2}
\lim_{N\to\infty}N^{2q+1}\|\widetilde R_{q,N}(f)\|^2
=\frac{2^{2q}}{\pi^{2q+2}}\bigl[(A+B)^2\,J_E+(A-B)^2\,J_O\bigr],
\end{equation}
where $A,B$ are the Dirac jump constants \eqref{eq:AB} and
\begin{equation}\label{eq:JEO}
J_E:=\int_{-1}^{1}\!\prod_{s\in E}\!\Bigl(x-\frac{1}{c_s}\Bigr)^{\!2}dx,
\qquad
J_O:=\int_{-1}^{1}\!\prod_{s\in O}\!\Bigl(x-\frac{1}{c_s}\Bigr)^{\!2}dx.
\end{equation}
\end{theorem}

\begin{proof} 
Since $\widetilde F_n$ is the generalized Fourier coefficient of
$f-\widetilde P$, we have
\[
\widetilde R_{q,N}(f)
=
\sum_{|n|>N}\widetilde F_n v_n.
\]
Hence, by Parseval,
\[
\|\widetilde R_{q,N}(f)\|^2
=
\sum_{|n|>N}|\widetilde F_n|^2.
\]
Since
\[
\widetilde F_n
=
F_n+(P_n-\widetilde P_n),
\]
we first compute the second term. By the definitions of $P_n$ and
$\widetilde P_n$,
\begin{align*}
P_n-\widetilde P_n
&=
-x_n\bigl[
\kappa_n^{(+)}\Delta^+(x_n)
-\kappa_n^{(-)}\Delta^-(x_n)
\bigr]
\\
&=
-x_n\bigl[
\kappa_n^{(+)}\Phi^+(x_n)
-\kappa_n^{(-)}\Phi^-(x_n)
\bigr]
\\
&\quad
-x_n^{q+1}
\bigl[
\kappa_n^{(+)}g_q^+
-\kappa_n^{(-)}g_q^-
\bigr],
\end{align*}
where we used \eqref{eq:Phi-def}. Define
\[
\mathcal Q_n
:=
\kappa_n^{(+)}\Phi^+(x_n)
-\kappa_n^{(-)}\Phi^-(x_n).
\]
Then \eqref{eq:Fn-asymp} gives
\[
P_n-\widetilde P_n
=
-x_n\mathcal Q_n-F_n+r_{n},
\qquad
|n|>N,
\]
where
\[
r_{n}=o(|n|^{-q-1})
\]
uniformly for $|n|>N$. Therefore
\begin{equation}\label{eq:tildeFn-id}
\widetilde F_n
=
-x_n\mathcal Q_n+r_{n}.
\end{equation}

By \eqref{eq:Phi-decomp-2},
\begin{align*}
\mathcal Q_n
&=
-g_q^+
\bigl[
\kappa_n^{(+)}U^+(x_n)
-\kappa_n^{(-)}U^-(x_n)
\bigr]
\\
&\quad
-g_q^-
\bigl[
\kappa_n^{(+)}V^+(x_n)
-\kappa_n^{(-)}V^-(x_n)
\bigr]
\\
&\quad
+\kappa_n^{(+)}\rho^+(x_n)
-\kappa_n^{(-)}\rho^-(x_n).
\end{align*}
Since
\[
\rho^\pm(x)
=
\sum_{k=0}^{q-1}\rho_k^\pm x^k,
\qquad
\rho_k^\pm=o(N^{k-q}),
\]
and $|x_n|=O(N^{-1})$ uniformly for $|n|>N$, it follows that
\[
\rho^\pm(x_n)=o(N^{-q}),
\]
uniformly for $|n|>N$ as $N\to\infty$.

Using \eqref{eq:kappa-leading} and \eqref{eq:UV-explicit}, we now
distinguish the parity of the tail index $n$. For $|n|>N$, we have
\[
U^\pm(x_n),V^\pm(x_n)=O(N^{-q}).
\]
Hence, if $n$ is even,
\[
\kappa_n^{(+)}U^+(x_n)-\kappa_n^{(-)}U^-(x_n)
=
K_0\Pi_E(x_n)
+
o\!\left(N^{-q}\right),
\]
and
\[
\kappa_n^{(+)}V^+(x_n)-\kappa_n^{(-)}V^-(x_n)
=
-K_0\Pi_E(x_n)
+
o\!\left(N^{-q}\right).
\]
If $n$ is odd,
\[
\kappa_n^{(+)}U^+(x_n)-\kappa_n^{(-)}U^-(x_n)
=
-K_0\Pi_O(x_n)
+
o\!\left(N^{-q}\right),
\]
and
\[
\kappa_n^{(+)}V^+(x_n)-\kappa_n^{(-)}V^-(x_n)
=
-K_0\Pi_O(x_n)
+
o\!\left(N^{-q}\right).
\]
Hence
\begin{equation}\label{eq:Qn-cases}
\mathcal Q_n
=
\begin{cases}
-K_0(g_q^+-g_q^-)\Pi_E(x_n)+o\!\left(N^{-q}\right),
& n\text{ even},
\\[2mm]
\phantom{-}K_0(g_q^++g_q^-)\Pi_O(x_n)+o\!\left(N^{-q}\right),
& n\text{ odd}.
\end{cases}
\end{equation}

Combining \eqref{eq:tildeFn-id} and \eqref{eq:Qn-cases}, and using
\eqref{eq:Pi-EO}, we obtain
\begin{equation}\label{eq:tildeFn-prod}
\widetilde F_n
=
\begin{cases}
\displaystyle
K_0(g_q^+-g_q^-)\,
x_n^{q+1}
\prod_{s\in E}
\left(1-\frac{n}{n_s}\right)
+x_n\sigma_{n,N},
& n\text{ even},
\\[4mm]
\displaystyle
-K_0(g_q^++g_q^-)\,
x_n^{q+1}
\prod_{s\in O}
\left(1-\frac{n}{n_s}\right)
+x_n\sigma_{n,N},
& n\text{ odd},
\end{cases}
\end{equation}
where
\[
\sup_{|n|>N}|\sigma_{n,N}|=o(N^{-q}).
\]
Here we used \eqref{eq:Pi-EO} and \eqref{eq:lam-asymp}. Indeed,
\[
nx_n=\frac{2}{\pi}+O\!\left(\frac{1}{n}\right),
\]
and hence, uniformly for $|n|>N$,
\[
(x_n-x_{n_s})
-
x_n\left(1-\frac{n}{n_s}\right)
=
\frac{nx_n-n_sx_{n_s}}{n_s}
=
O(N^{-2}).
\]
It follows that
\[
x_n\Pi_E(x_n)
=
x_n^{q+1}
\prod_{s\in E}
\left(1-\frac{n}{n_s}\right)
+
x_nO(N^{-q-1}),
\]
and analogously for $\Pi_O$.

The remainders in \eqref{eq:tildeFn-prod} are of the form
$x_n\sigma_{n,N}$, where
\[
\sup_{|n|>N}|\sigma_{n,N}|=o(N^{-q}).
\]
Since
\[
\sum_{|n|>N}|x_n|^2=O(N^{-1}),
\]
we have
\[
\sum_{|n|>N}|x_n\sigma_{n,N}|^2
=
o(N^{-2q-1}).
\]
The squared sum of the leading terms in \eqref{eq:tildeFn-prod} is
$O(N^{-2q-1})$, so the corresponding cross terms are also
$o(N^{-2q-1})$ by the Cauchy--Schwarz inequality. Thus only the
leading terms in \eqref{eq:tildeFn-prod} contribute to the limit.

For the even indices, the Riemann-sum argument gives
\[
\begin{aligned}
N^{2q+1}
\sum_{\substack{|n|>N\\ n\ {\rm even}}}
|\widetilde F_n|^2
\longrightarrow
K_0^2(g_q^+-g_q^-)^2
\frac{(2/\pi)^{2q+2}}{2}
\int_{|u|>1}
u^{-2q-2}
\prod_{s\in E}
\left(1-\frac{u}{c_s}\right)^2du.
\end{aligned}
\]
The factor $1/2$ accounts for the density of the even integers.
With the change of variables $u=1/x$,
\[
\int_{|u|>1}
u^{-2q-2}
\prod_{s\in E}
\left(1-\frac{u}{c_s}\right)^2du
=
\int_{-1}^{1}
\prod_{s\in E}
\left(x-\frac{1}{c_s}\right)^2dx
=
J_E.
\]
Similarly,
\[
N^{2q+1}
\sum_{\substack{|n|>N\\ n\ {\rm odd}}}
|\widetilde F_n|^2
\longrightarrow
K_0^2(g_q^++g_q^-)^2
\frac{(2/\pi)^{2q+2}}{2}
J_O.
\]

Using $K_0^2=1/2$, we conclude that
\[
N^{2q+1}\|\widetilde R_{q,N}(f)\|^2
\longrightarrow
\frac{2^{2q}}{\pi^{2q+2}}
\left[
(g_q^+-g_q^-)^2J_E
+
(g_q^++g_q^-)^2J_O
\right].
\]
Finally, by \eqref{eq:AB},
\[
g_q^+-g_q^-=A+B,
\qquad
g_q^++g_q^-=A-B,
\]
so this is precisely \eqref{eq:L2}.
\end{proof}

\end{document}
